\chapter{Feedback Loops and Digital Twins}

\section{Entering Part V}

The previous parts of this book have treated constraint projection, the active correction that keeps the reactor within its admissible region, largely as a given: something that happens, performed by unspecified control systems and unspecified operators, without asking in detail how it actually happens or what it depends on to keep happening. Part V opens that question up. It asks how feedback loops, both automated and human, actually sustain constraint projection across the reactor's operating life, and it treats the answer as recursive in a specific sense: the mechanisms that stabilize the reactor are themselves systems that must be stabilized, by further mechanisms, in a nested structure that does not terminate in any single, self-sufficient controller sitting outside the loop it manages.

\section{Feedback as the Physical Mechanism of Projection}

Every instance of constraint projection described so far, a plasma control system suppressing a growing instability, a cooling system responding to a thermal transient, an operator adjusting a heating profile, is implemented through feedback: a sensor measures some aspect of the reactor's current state, that measurement is compared against a target or an acceptable range, and the discrepancy drives a correction intended to reduce it. This is true whether the loop closes in microseconds through dedicated control electronics or over months through an operator's accumulated judgment about how the plant tends to drift. Feedback is not one particular technology among several available for reactor control. It is the general form that constraint projection takes, at every timescale, whenever a system's own past behavior is used to shape its future behavior toward some target.

This has an immediate consequence for how feedback systems should be evaluated: not merely by how accurately they measure the reactor's state, but by how well their correction, once applied, actually moves the state back toward the interior of the admissible region rather than overcorrecting past it or responding too slowly to prevent a boundary crossing. A sensor system that measures plasma density with excellent precision but whose correction signal reaches the confinement field too slowly to counteract a fast-growing instability provides little practical value, however accurate its measurement, because the value of feedback lies entirely in the closed loop it forms, not in the isolated quality of the measurement it starts from.

\section{Layered Feedback}

No single feedback loop handles the entire span of timescales this book has discussed, and it would be a design error to expect one to. Instead, real fusion reactors, like most complex engineered systems that must remain viable over long periods, rely on layered feedback: fast, narrowly scoped loops nested inside slower, more broadly scoped ones, each layer correcting deviations its own timescale is suited to detect and address, and each layer's stable operation providing the relatively unchanging background against which the next slower layer can operate effectively.

A plasma stability control loop, operating on millisecond timescales, corrects instabilities without needing to know anything about the reactor's tritium inventory or the condition of its structural materials. An operational planning loop, operating on timescales of hours to days, adjusts heating profiles, cooling flow rates, and extraction rates in response to trends the fast loop does not track, without needing to model the plasma's instantaneous behavior in detail. A maintenance planning loop, operating on timescales of months to years, schedules component replacement and refurbishment in response to accumulated degradation the operational loop does not track, without needing to model day-to-day operational variation directly. Each layer depends on the layer below it functioning well enough that the assumptions it relies on, essentially fixed conditions at its own timescale, remain approximately true. This is the same timescale separation discussed in Chapter 8, now viewed specifically through the lens of what makes it possible: each layer is itself a feedback loop, and the layered structure as a whole is a feedback loop of feedback loops.

\section{The Recursive Structure of Stabilization}

This layering reveals something that a simpler picture of reactor control tends to obscure: there is no single controller, sitting entirely outside the system it manages, that stabilizes the reactor as a whole. The fast plasma control loop is stabilized, in the sense of having its own assumptions kept valid, by the slower operational loop maintaining conditions within the range the fast loop was designed for. The operational loop, in turn, is stabilized by the maintenance loop keeping components within a condition the operational loop's models assume. The maintenance loop is stabilized by the institutional knowledge loop, discussed in the previous chapter, ensuring that the people planning and executing maintenance retain the understanding needed to do so correctly. And the institutional knowledge loop is itself stabilized, in part, by the economic and regulatory conditions taken up later in this Part, which determine whether adequate staffing, training, and investment in maintenance capacity remain available at all.

Each layer, in other words, is not stabilized from outside the recursive structure but by the layer adjacent to it within that same structure, which is itself stabilized by the layer adjacent to it, and so on. This is why this book has consistently resisted describing the reactor as a machine with an external operator standing apart from it, making decisions about a passive object. The more accurate picture is a nested set of feedback processes, each dependent on the others remaining functional, together constituting the reactor's overall viability, with no privileged layer that could be removed while leaving the rest intact.

\section{Digital Twins as Prediction Folded Into Repair}

A digital twin, a continuously updated computational model of the reactor that mirrors its current state and can simulate its likely future behavior under various conditions, fits naturally into this picture, though it is worth being precise about what role it actually plays. A digital twin does not replace any of the feedback layers described above. It augments them, by allowing a given layer to anticipate a deviation before it has fully manifested in the physical system, rather than reacting only after a sensor has registered the deviation directly.

This is, in the language this book has developed, prediction functioning as a form of continual repair, extending the effective response time available to a given feedback layer by acting on a forecast rather than waiting for the disturbance itself to arrive. A digital twin that accurately models how a particular pattern of plasma behavior tends to precede a disruption allows the fast control layer to intervene earlier, with a gentler correction, than it could if it were only responding to the disruption's own early physical signs. A digital twin that accurately models how a given operating history tends to accelerate structural degradation allows the maintenance layer to schedule intervention before degradation reaches a level that would otherwise only be discovered during a routine inspection, after some of the available margin has already been consumed.

The value of a digital twin, in this framing, is precisely proportional to the accuracy of its predictions and the speed with which those predictions can be turned into action within the relevant feedback layer. A highly accurate model whose predictions cannot be acted on quickly enough provides little practical benefit over waiting for direct measurement, and a model that can be acted on quickly but whose predictions are unreliable risks introducing corrections based on forecasts that do not actually describe what the reactor is about to do, potentially destabilizing a layer that would have performed better relying on direct measurement alone. A digital twin is therefore not simply a tool to be added to the reactor's control architecture wherever computational resources permit. It is a component whose value depends on how well it is matched to the specific feedback layer it is meant to augment, in the same way that every other repair and stabilization mechanism discussed in this book has depended on being matched to the specific timescale and failure mode it addresses.

\section{Toward the Human and Economic Layers}

This chapter has described the reactor's feedback structure largely in terms of control systems, sensors, and computational models, but it has already gestured toward the layers that lie beyond pure automation: the operational judgment that interprets ambiguous conditions no automated system was designed for, the institutional knowledge that determines whether maintenance planning is executed correctly, and the economic and regulatory conditions that determine whether adequate resources remain available to sustain any of this at all. The next two chapters take up these layers directly, treating human operators and supply chains, and then economic viability itself, as further instances of the same recursive feedback structure developed here, rather than as separate concerns external to the reactor's technical viability.
